\documentclass[conference]{IEEEtran}
\usepackage{times}

% numbers option provides compact numerical references in the text. 
\usepackage[numbers]{natbib}
\usepackage{multicol}
\usepackage[bookmarks=true]{hyperref}
\usepackage{amssymb}
\usepackage{amsmath}
\usepackage{booktabs}


\pdfinfo{
   /Author (Homer Simpson)
   /Title  (Robots: Our new overlords)
   /CreationDate (D:20101201120000)
   /Subject (Robots)
   /Keywords (Robots;Overlords)
}

\begin{document}

% paper title
\title{EmoMoconvq: Emotion Conditioned Human Motion Generation}

% You will get a Paper-ID when submitting a pdf file to the conference system
\author{Kangwei Sun 524531910006}

%\author{\authorblockN{Michael Shell}
%\authorblockA{School of Electrical and\\Computer Engineering\\
%Georgia Institute of Technology\\
%Atlanta, Georgia 30332--0250\\
%Email: mshell@ece.gatech.edu}
%\and
%\authorblockN{Homer Simpson}
%\authorblockA{Twentieth Century Fox\\
%Springfield, USA\\
%Email: homer@thesimpsons.com}
%\and
%\authorblockN{James Kirk\\ and Montgomery Scott}
%\authorblockA{Starfleet Academy\\
%San Francisco, California 96678-2391\\
%Telephone: (800) 555--1212\\
%Fax: (888) 555--1212}}


% avoiding spaces at the end of the author lines is not a problem with
% conference papers because we don't use \thanks or \IEEEmembership


% for over three affiliations, or if they all won't fit within the width
% of the page, use this alternative format:
% 
%\author{\authorblockN{Michael Shell\authorrefmark{1},
%Homer Simpson\authorrefmark{2},
%James Kirk\authorrefmark{3}, 
%Montgomery Scott\authorrefmark{3} and
%Eldon Tyrell\authorrefmark{4}}
%\authorblockA{\authorrefmark{1}School of Electrical and Computer Engineering\\
%Georgia Institute of Technology,
%Atlanta, Georgia 30332--0250\\ Email: mshell@ece.gatech.edu}
%\authorblockA{\authorrefmark{2}Twentieth Century Fox, Springfield, USA\\
%Email: homer@thesimpsons.com}
%\authorblockA{\authorrefmark{3}Starfleet Academy, San Francisco, California 96678-2391\\
%Telephone: (800) 555--1212, Fax: (888) 555--1212}
%\authorblockA{\authorrefmark{4}Tyrell Inc., 123 Replicant Street, Los Angeles, California 90210--4321}}


\maketitle

\begin{abstract}
Physics-based motion generation methods such as MoConVQ~\cite{yao2023moconvq}
produce high-quality motions from text descriptions but generate stylistically neutral
results regardless of emotional context. We propose \textbf{EmoMoconvq}, which extends
MoConVQ's T2M-MoConGPT by injecting a discrete emotion label
(e.g., \textit{happy}, \textit{sad}, \textit{angry}, \textit{fearful}) as an
additional conditioning signal. Through parameter-efficient fine-tuning on a small
emotionally annotated motion dataset, we aim to generate physically plausible motions
whose stylistic qualities reflect the intended emotional state.
\end{abstract}

\IEEEpeerreviewmaketitle

\section{Motivation and Problem Statement}
Human motion is inherently expressive---the same action performed in different emotional
states (e.g., walking while happy vs.\ sad) exhibits distinct kinematic patterns such
as stride length, torso posture, and arm swing amplitude. This expressiveness is critical
for digital humans, game characters, and social robots, yet current physics-based
generation methods ignore it entirely.

MoConVQ~\cite{yao2023moconvq} shows that discrete VQ-VAE representations combined
with a GPT-style transformer can generate diverse, physically accurate motions from
natural language. However, its T2M-MoConGPT produces emotion-neutral outputs: the same
text prompt always yields the same stylistic result. We therefore propose \textbf{EmoMoconvq},
which extends T2M-MoConGPT by conditioning on discrete emotion labels. 

\section{Literature Review}

\subsection{Physics-Based Motion Generation.}
MoConVQ~\cite{yao2023moconvq} encodes motions into residual VQ codes and decodes
them through a physics simulator via model-based RL, trained on 23 hours of diverse
motion data. Its frozen encoder and codebook make it well-suited for lightweight
downstream fine-tuning.

\subsection{Text-to-Motion Generation.}
T2M-GPT~\cite{zhang2023t2m} and MDM~\cite{tevet2022human} generate kinematic motions
from text using discrete transformers and diffusion models respectively, both trained
on HumanML3D~\cite{guo2022generating}. Neither is physics-based nor models emotional style.

\subsection{Style and Emotion in Motion.}
ACTOR~\cite{petrovich2021action} conditions a Transformer VAE on action labels, showing
that high-level categorical signals can guide motion style. The CMU MoCap
database~\cite{cmu_mocap} includes emotionally performed locomotion sequences suitable
for fine-tuning. Our work is the first to combine emotion conditioning with a
physics-based generation pipeline.



\section{Technical Plan}

\subsection{System Overview}
 
We build directly on the pretrained MoConVQ framework~\cite{yao2023moconvq}.
The MoConVQ system consists of three trained components: (1) a 1D convolutional
\textbf{motion encoder} $\mathcal{E}$ that maps a motion clip $M$ to a sequence of
latent vectors $V = [v^k]$; (2) a \textbf{residual VQ codebook} $\mathcal{B}$ with
$N=8$ quantization layers, each containing 512 code vectors of dimension 768, producing
discrete index sequences $S = [I^k]$; and (3) a \textbf{physics-based decoder}
$\mathcal{D}$ comprising a deconvolutional network, a mixture-of-experts control policy,
and a physics simulator. For text-to-motion generation, MoConVQ employs a dual-transformer
\textbf{MoConGPT}: a temporal transformer ($N_T = 12$ layers) that aggregates information
across motion time steps, and a depth transformer ($N_D = 5$ layers) that predicts
RVQ indices layer by layer within each time step.
 
In our extension, we freeze all components except for a newly introduced
\textbf{emotion embedding module} and the \textbf{top $K$ layers} of the temporal
transformer.

\subsection{Emotion Embedding Module}
 
Let $\mathcal{C} = \{\textit{neutral},\, \textit{happy},\, \textit{sad},\,
\textit{angry},\, \textit{fearful}\}$ denote the set of $|\mathcal{C}|=5$ emotion
categories. Each emotion $c \in \mathcal{C}$ is mapped to a trainable embedding vector:
\begin{equation}
    E_e: \mathcal{C} \rightarrow \mathbb{R}^{d_e}, \quad e_c = E_e(c)
\end{equation}
where $d_e = 512$ matches the hidden dimension of the temporal transformer. The
emotion embedding is fused with the T5 text feature $f_{\text{T5}}$ via a two-layer
MLP projection followed by element-wise addition:
\begin{equation}
    f_{\text{cond}} = f_{\text{T5}} + \text{MLP}_{\text{proj}}(e_c)
\end{equation}
This additive fusion preserves the pretrained text conditioning while introducing an
orthogonal affective signal. The fused feature $f_{\text{cond}}$ replaces $f_{\text{T5}}$
in the cross-attention layers of the temporal transformer, following the same injection
mechanism as the original T2M-MoConGPT.
 
\subsection{Fine-Tuning Strategy}
 
We freeze the MoConVQ encoder, codebook, physics decoder, depth transformer, and the
bottom $(N_T - K)$ layers of the temporal transformer. We fine-tune only:
\begin{itemize}[noitemsep]
    \item The emotion embedding table $E_e$ ($|\mathcal{C}| \times d_e$ parameters),
    \item The MLP projection head $\text{MLP}_{\text{proj}}$,
    \item The top $K = 4$ layers of the temporal transformer (cross-attention and FFN weights).
\end{itemize}
This parameter-efficient design minimizes catastrophic forgetting of the pretrained
motion distribution while allowing the model to learn emotion-specific stylistic
variations. We train with the same negative log-likelihood loss as MoConGPT:
\begin{equation}
    \mathcal{L}_{\text{emo}} = \mathbb{E}_{(S,\, c)} \left[ \sum_{k=1}^{K} \sum_{d=0}^{N}
    -\log p\!\left(I^k_d \,\middle|\, f_{\text{cond}}(c),\; I^k_{<d},\; I^{<k}\right) \right]
\end{equation}
  
\subsection{Data Preparation}
 
We construct a fine-tuning dataset by combining two sources:
 
\textbf{(1) HumanML3D subset with emotion keywords.} We filter the HumanML3D training
set~\cite{guo2022generating} for text annotations containing explicit emotion-related
adverbs or adjectives (e.g., ``happily'', ``angrily'', ``fearfully'', ``sadly'').
We assign emotion labels based on keyword matching and manual verification. We expect
approximately 800--1200 clip-text pairs across the four non-neutral emotion classes.
 
\textbf{(2) CMU MoCap emotional gait sequences.} The CMU database includes a small
set of locomotion sequences performed in specified emotional styles. We retarget these
sequences to the MoConVQ character skeleton and synthesize paired text descriptions
(e.g., ``a person walks forward'') with assigned emotion labels. This yields
approximately 100--200 additional pairs.
 
For data augmentation, we apply random horizontal mirroring (left/right reflection)
and temporally subsample sequences at 0.8$\times$ and 1.2$\times$ speed, doubling
the effective dataset size.
 
\subsection{Evaluation Protocol}
 
We evaluate on a held-out set of emotion-conditioned text prompts and report:
 
\begin{itemize}[noitemsep]
    \item \textbf{FID}: Fr\'{e}chet Inception Distance between generated and
    ground-truth motion feature distributions, computed using the pretrained motion
    feature extractor from~\cite{guo2022generating}.
    \item \textbf{R-Precision (Top-1/2/3)}: Retrieval accuracy between generated
    motions and their conditioning text descriptions.
    \item \textbf{Emotion Classification Accuracy}: A separately trained lightweight
    emotion classifier (e.g., GRU-based) on motion sequences is used to verify that
    generated motions are classified as the intended emotion.
    \item \textbf{User Study}: 20 participants rate generated motion pairs
    (same action, different emotion labels) on a 5-point Likert scale for
    (a) emotional distinctiveness and (b) motion naturalness.
\end{itemize}

\section{Expected Outcomes}
 
We expect to demonstrate that EmoMotion successfully conditions motion style on
discrete emotion labels without sacrificing the physics-based plausibility guaranteed
by the MoConVQ decoder. Concretely, we anticipate:
\begin{itemize}[noitemsep]
    \item FID scores competitive with or modestly higher than the baseline
    T2M-MoConGPT (target: $\leq 0.5$ FID degradation),
    \item Emotion classification accuracy $\geq 60\%$ on held-out clips
    (significantly above the $20\%$ random baseline for 5 classes),
    \item User study ratings indicating that $\geq 70\%$ of participants can
    correctly identify the intended emotion from the generated motion,
    \item Qualitative video results showing perceptibly different walking,
    running, and sitting motions across emotion conditions.
\end{itemize}

\section{Potential Risks and Backup Plan}

\subsection{Insufficient labeled data.} 
If fine-tuning overfits, we reduce to 3 emotion
classes (happy/sad/neutral) and rely more heavily on augmentation.

\subsection{Catastrophic forgetting.} 
If overall FID degrades sharply, we reduce $K$
(freeze more layers) or adopt LoRA adapters~\cite{hu2022lora} instead of full
layer fine-tuning.

\subsection{Compute constraints.} As a fallback, we fine-tune only the emotion embedding
and MLP ($<$1M parameters), omitting temporal transformer fine-tuning.

\section{Timeline}

\begin{center}
\small
\begin{tabular}{p{2.5cm} p{5cm}}
\toprule
\textbf{Period} & \textbf{Milestone} \\
\midrule
Week 1--2  & Reproduce T2M-MoConGPT baseline; verify FID/R-precision. \\
Week 3--4  & Curate emotion dataset; implement augmentation pipeline. \\
Week 5--6  & Implement EmoMoconvq modules; run initial fine-tuning. \\
Week 7--8 & Ablation studies; quantitative evaluation; user study. \\
Week 9--10 & Final report and demo video. \\
\bottomrule
\end{tabular}
\end{center}



%% Use plainnat to work nicely with natbib. 

\bibliographystyle{plainnat}
\bibliography{references}

\end{document}


